\part{2025}
\section{
    0926联合早报·下午察：被清朗的户晨风和张雪峰
} {
    \setlength{\parindent}{2em}
    中国互联网本周可谓风声鹤唳。无论是凭借“安卓人”“苹果人”等网络热梗爆红的户晨风，还是高调喊话要在大陆攻台时捐出1亿元（人民币，下同，1801万新元）的考研名师张雪峰，都在多个平台遭遇“禁声”。

    舆论一时哗然，外界普遍猜测，这一连串动作并非巧合，而是中国新一轮整治网络负面情绪行动的一环。

中共中央网信办9月22日在官网公告，将在全国部署开展为期两个月的“清朗·整治恶意挑动负面情绪问题”专项行动，以营造更加文明理性的网络环境。

公告显示，此次专项行动聚焦社交、短视频、直播等平台，全面排查话题、榜单、推荐、弹幕、评论等重点环节，着力整治挑动群体极端对立情绪、宣扬恐慌焦虑情绪、挑起网络暴力戾气，以及过度渲染消极悲观情绪等问题。

这波操作并非无的放矢。近年来，中国经济因房地产危机陷入低迷，消费需求被抑制，就业压力也随之加剧。面对失业率居高不下，以及高校与岗位间的激烈竞争，许多年轻人已对社会和经济前景生出幻灭感。

新加坡南洋理工大学社会科学学院助理教授骆斯航对BBC说，中国年轻人“对未来生活的前景存有严重疑问”，而且“必须直面一个现实，即他们的生活很可能比父辈更差”。

随着这种不满情绪持续累积，北京开始重拳出击，而第一波整治对象，正是网络红人。

“苹果人、安卓人”
第一个“中招”的就是户晨风。9月20日，他在微博、B站、抖音、快手等多个平台的账号被禁言，所有内容清空，理由是“违反相关法律法规”。

而在“消失”之前，资深“果粉”户晨风在直播时的一句“你是典型的安卓逻辑、安卓人、安卓学历、安卓素质、安卓态度”，意外引爆全网；“安卓人”“苹果人”很快刷屏，成为近期最火的流行梗之一。

《广州日报》评论认为，户晨风被禁言，就是因为他用“苹果”“安卓”标签，把消费、学历和手机品牌强行挂钩，制造出一条“鄙视链”。

在户晨风的叙事里，“安卓人”意味着底层、低收入、没素质，“安卓学历”指的是专科生等较低学历；而“苹果人”则是体面、高端、收入更高的象征。

这套“苹果人、安卓人”逻辑走红后，户晨风还顺势推出“祝你越来越苹果”的短视频业务，几十秒的视频标价2999元，却照样有人买单。他还在直播中透露，自己的月收入已经突破60万元。

但随着热度攀升，户晨风的言论也被指变相挑动阶级对立、传播负面情绪，正好踩中了这轮整治的红线——网信办明确点名要整治“炮制热梗……过度自我矮化或者渲染颓丧或消极情绪”的行为。

户晨风被禁言的消息一出，网上炸开了锅。有人拍手叫好，直呼“活该”；但也有人觉得这些梗只是个笑话，不必小题大做。

也有网民指出，户晨风本身只有高中学历，是靠拍关怀底层的视频一步步走上来的。但随着人气飙升，他却开始嘲讽“安卓人”，似乎忘了自己也曾是“安卓人”。

有趣的是，就在户晨风被封的同时，苹果新机开售。在中国大陆乃至香港、日本、新加坡、美国等全球多地的苹果门店里，不少人把展示机壁纸换成了户晨风的照片。这场“恶搞”，反而让这位“被清朗”的网红更出圈了。

在一些平台上，户晨风的口碑甚至开始逆转。抖音上一条高赞评论写道：“告诉户子，无论他在哪个平台重生，我都会去迎接他。”

事实上，这并不是户晨风首次被封。自2023年进军短视频以来，他因拍摄“百元养老金购买力挑战”，以及在直播连线中触及政治敏感话题，多次被短暂禁言。

微博大V、中国体育记者王玮晨直言，户晨风“吃的都是危险流量，怂恿阶级对立、贬低国产产品是大忌，不可能不办你”。

知名网红“留几手”也说：“户子总觉得注意敏感词就行了，可他有没有想过，把一群人聚在一起瞎折腾，解构行业、解构阶层，把很多严肃的东西戏谑化，本身就是敏感。”

张雪峰“捐1亿”也难逃禁声？
另一名“中招”的，是考研名师张雪峰。9月24日，他在微博、小红书、抖音、B站等平台的账号均被禁止新增关注。

张雪峰团队解释，这是因为他在直播中出现不当言论，违反平台规则，遭到举报，“公司已经在处理，应该短期内就可以解决”，并称“正在反省”。但团队并未说明具体是哪场直播、哪句话触发封禁。

外界普遍将此与张雪峰9月3日阅兵后的表态联系起来。当时，张雪峰强调中国武器“不是为了吓唬和干别人，而是为了国家统一”，并多次爆粗口称，如果大陆攻台，他至少会捐出1亿元。

相关视频曝光后，在网上引发争议。张雪峰在直播中回应此事时说：“爱咋咋地，老子不移民。我怕我就不会说，说了我就一定会做。”

张雪峰当时还称，当地网信办官员甚至提醒他“注意一下”，视频被发到了境外网络平台；他随即直呼“群众当中有坏人”。

表面上，张雪峰的立场与中国官方的“祖国统一”论调一致，但话里涉及战争与焦虑，正好踩到这轮清朗行动的雷区。

其实，今年41岁的张雪峰本身就非常具有争议性。他曾在为学生和家长提供填报大学志愿建议的直播中称：“所有文科专业都叫服务业，服务业总结成一个字就是‘舔’。”他还极力不建议选择新闻专业，直言如果孩子坚持报考，“一定会把他打晕”。

而网信办的公告中提到，此次整治行动将重点打击“编造‘大师’‘专家’虚假身份人设，围绕就业、婚恋、教育等领域贩卖焦虑带货卖课”的现象。在不少网民看来，张雪峰虽然打着帮学生指点迷津的旗号，但言论常常制造焦虑，正好撞上了这次清朗行动的枪口。

中国网络的边界
实际上，这并非清朗行动首次出现。自2016年以来，网信办就不断以“清朗”为名开展专项整治，范围涵盖网络暴力、假消息以及恶意营销等，意在打造“天朗气清、生态良好”的网络环境。

长期以来，中国对网络内容实行严格管控。虽然西方平台也会规范用户行为，但中国的管控范围更广，背后的考量是担心网上的激烈争论和情绪煽动会影响现实社会秩序。

近期研究显示，中国民众对未来前景的悲观情绪正不断加剧。专家指出，中国官方显然注意到了这一点，并试图通过整治行动压制相关迹象，但效果如何，仍不得而知。

骆斯航认为：“表达悲观情绪，并不等于否定参与劳动力市场或整个社会的可能性。但如果剥夺了人们宣泄情绪后的解脱感，反而可能让群体心理状况更糟。”

澎湃新闻则发表评论文章称，这次清朗行动“正当其时”，既是对当前乱象的回应，也为网络空间的运转划定底线。

无论是户晨风还是张雪峰，他们的账号被封，都在言论出圈、广受热议之后，可谓”正当其时”。这似乎说明，本轮清朗行动瞄准的不只是负面情绪，更是任何可能搅动社会不稳定的因素。

只是，即便网络环境能被快速“清理”，人们心中真实的不满与悲观情绪，恐怕不是那么轻易地挥之而去。

source: https://www.zaobao.com/realtime/china/story20250926-7574995?ref=home-top-news
}